\title{LongRoPE: Extending LLM Context Window Beyond 2 Million Tokens}

\begin{abstract}
A substantial context window is an advantageous characteristic for large language models (LLMs). Nonetheless, existing context window extensions are generally restricted to approximately 128k tokens. This limitation arises from the high costs associated with fine-tuning, the limited availability of lengthy textual data, and the disruptive effects of new token positions that can introduce catastrophic values.
\end{abstract}

In this work, we present LongRoPE, a novel approach that successfully extends the context window of existing pre-trained LLMs to a remarkable \textbf{2048k} tokens. This expansion is accomplished with as few as 1k fine-tuning steps, conducted at training lengths up to 256k, all while preserving the original model’s performance for shorter context lengths. Our approach relies on three major contributions: (i) we systematically detect and utilize two distinct types of non-uniformities in positional interpolation through an efficient search, which offers a superior starting point for fine-tuning and allows for an 8$\times$ window extension even without additional fine-tuning; (ii) we propose a staged extension technique, initially fine-tuning the LLM at a 256k context size, followed by a secondary positional interpolation on the already-extended model to reach a 2048k context capacity; (iii) we further recalibrate LongRoPE at 8k length, restoring and maintaining high-level performance for short context usage. Comprehensive evaluations using LLaMA2 and Mistral across a variety of benchmarks validate the effectiveness of our technique. Models adapted through LongRoPE require only minimal changes to the positional embedding, retaining the original architecture and allowing continued use of most established optimizations. The implementation will be released at \texttt{https://github.com/microsoft/LongRoPE}

\section{Introduction}

Although Large Language Models (LLMs) have achieved significant advancements across a wide range of applications~\cite{gpt4,llama2}, they are constrained by relatively short context windows, such as the 4096-token threshold in LLaMA2~\cite{llama2}. When the input exceeds this predefined length, model performance degrades, primarily because these positions were not encountered during training. This limitation presents obstacles for use cases that require processing extensive in-context examples~\cite{Huang2023FewerIM} or for powering LLM-based agents~\cite{park2023generative,madaan2023selfrefine}.

Recent studies have demonstrated that the context window of a pre-trained LLM can be enlarged to approximately 128k by employing fine-tuning strategies on longer textual sequences~\cite{longlora,pi,yarn,zhang2024soaring,liu2023scaling}. Despite these advancements, three principal challenges impede the expansion of the context window beyond this range. Firstly, introducing novel, untrained positional indices produces a multitude of problematic, outlier values, exacerbating distributional shifts and hindering convergence during fine-tuning~\cite{pi}. This issue becomes pronounced in scenarios where the jump is significant—for example, increasing the context length from 4k to more than 1000k involves introducing over 90\% entirely new positional embeddings. Secondly, successful fine-tuning typically depends on the availability of text samples that match these extended lengths; yet, texts exceeding 1000k are scarce in existing datasets. Furthermore, processing such lengthy documents is computationally intensive, often making the training process excessively time-consuming and resource-demanding. Thirdly, as the context window reaches extreme lengths, attention scores are allocated across a vast number of token positions, diluting focus and thereby impairing performance on the model's originally intended, shorter sequences~\cite{pi}.

A prevalent strategy for addressing the initial issue involves interpolating RoPE positional embeddings~\cite{rope,pi}, a process that rescales novel position indices to the pretrained interval, as illustrated in Fig.\ref{fig:overview}. The Position Interpolation (PI) technique~\cite{pi} achieves this by linearly interpolating the rotary angles of RoPE according to the ratio of extension. In contrast, NTK~\cite{ntk,dynamicntk} recommends applying unequal interpolation and extrapolation along the RoPE dimensions. The YaRN approach~\cite{yarn} divides the RoPE dimensions into three groups based on their frequency, deploying extrapolation, NTK, and linear interpolation for each respective group. Nonetheless, positional embeddings within the Transformer architecture inherently possess \emph{complex, heterogeneous information entropy}. Present methods do not fully exploit this intricate non-uniformity, resulting in a loss of information that ultimately constrains the extendable context window.

Section~\ref{sec:analysis} identifies two principal empirical observations: \textbf{(1)} Successful positional interpolation requires accounting for two distinct sources of non-uniformity: the diversity across RoPE dimensions and across token positions. Interpolating less at smaller RoPE dimensions and for early token indices proves advantageous, although the ideal interpolation approach is contingent on the specific desired sequence length extension.
  \textbf{(2)} Incorporating these non-uniform characteristics into the positional interpolation process facilitates the preservation of information embedded in the original RoPE representation, particularly within key dimensions and at crucial token locations. This strategy reduces information degradation introduced by interpolation and, as a result, offers superior initialization for subsequent fine-tuning. Furthermore, this method enables sequence length extension by up to 8$\times$ even in the absence of fine-tuning.

Building upon these insights, we present \textbf{LongRoPE}, a robust approach designed to push the LLM context window past 2 \emph{million} tokens. LongRoPE leverages three principal innovations. To begin, {LongRoPE} takes full advantage of the diverse multidimensional characteristics found in positional interpolation. It determines optimal rescale factors for the rotation angles within RoPE, customized for each dimension and depending on the positions of tokens. Because the search space for these rescale factors increases exponentially as the desired extension ratio grows, {LongRoPE} incorporates an evolutionary search algorithm augmented by two optimization strategies to enhance search effectiveness. An illustration of a discovered rescaled RoPE is provided in Fig.~\ref{fig:overview}.

Subsequently, {LongRoPE} employs a progressive and efficient expansion approach to obtain a 2048k context window, circumventing the necessity of directly fine-tuning the model on ultra-long texts, which are rare and difficult to acquire. The approach initiates by exploring the performance at a 256k context length using the pre-trained LLM, followed by fine-tuning the model at this length. Leveraging our non-uniform positional interpolation method—which permits an 8$\times$ increase in the context window without further fine-tuning—a secondary search for optimal RoPE rescale factors is then performed on the extended and fine-tuned LLM. This two-stage strategy enables the extension to a 2048k context window on both LLaMA2 and Mistral~\cite{mistral}.

In order to prevent a drop in performance for shorter, original context windows, {LongRoPE} maintains its adjustment of the RoPE rescale parameters even after the LLM has been extended. Analogous to the upscaling process from 256k to 2048k, the method applies a downscaling strategy for context windows of 4k and 8k in the LLM that has been fine-tuned at 256k, utilizing the search algorithm to reduce the degree of positional interpolation. When processing input sequences shorter than 8k during inference, the RoPE is modified with the rescale factors identified through this search.

Comprehensive experimental results with multiple LLMs across a range of long-context tasks validate the superiority of our approach. Our findings indicate that LongRoPE consistently preserves low perplexity when evaluated on sequence lengths from 4k up to 2048k. It attains more than 90\% passkey retrieval accuracy and matches the performance of baseline models on established benchmarks built for a 4096 context window. Moreover, LongRoPE is universally compatible with any LLMs utilizing RoPE embeddings. We intend to make our codebase and the LongRoPE-2048k models publicly accessible.

\section{Non-uniformity in Positional Interpolation}
\label{sec:analysis}

\subsection{Preliminary}

Transformer architectures inherently depend on positional encoding—commonly known as position embedding—to capture the sequential arrangement of input tokens. In this study, we center our attention on the RoPE~\cite{rope} position embedding, a prevalent approach in contemporary LLMs. The RoPE representation assigned to a token located at position $n$ can be succinctly expressed as:

\begin{equation}
	\fontsize{7.8}{7.8} \selectfont
	\label{eq:rope}
	\left[ cos(n\theta_0), sin(n\theta_0), cos(n\theta_1), \cdot\cdot\cdot, cos(n\theta_{d/2-1}), sin(n\theta_{d/2-1}) \right]   
\end{equation}

Here, $d$ stands for the dimension of the embedding, and $n\theta_i$ denotes the angular value corresponding to the token’s $n$-th position. The rotation frequencies are defined by $\theta_i = \theta^{-2i/d}$. In the context of RoPE, the default value chosen for $\theta$ is 10000.

\textbf{\emph{Context window extension ratio $s$} and positional interpolation}. We define $s$ as  the ratio of extended context length $L'$ to the original length $L$: $s=\frac{L'}{L}$. 

In order to increase the context length from $L$ to $L'$, prevalent approaches using positional interpolation advocate reducing the rotation frequencies $\theta_i$ according to the context extension ratio $s$. Defining $\beta = \theta^{2/d}$ and letting $\lambda$ represent the effective scaling factor corresponding to $s$, we standardize these positional interpolation approaches in the following expression:

\begin{equation}	
	\fontsize{7.5}{7.5} \selectfont
	\label{eq:unified_rope}
	\begin{aligned}
		\left[cos\left(\frac{n}{\lambda(\beta)^0}\right), sin \left(\frac{n}{\lambda(\beta)^0}\right), cos\left(\frac{n}{\lambda(\beta)^1}\right), \cdots,  sin \left(\frac{n}{\lambda(\beta)^{d/2-1}}\right) \right]   
	\end{aligned}
\end{equation}

\textbf{Linear positional interpolation (PI)}. PI~\cite{pi} employs a linear scaling of position indices when working within the maximum pre-trained sequence length. When extending to a new sequence length with a target ratio $s$, the model compresses the rotation angles for every position by applying a factor of $\lambda=s$ across each RoPE dimension. This approach, however, leads to overly compressed positional encoding, resulting in dense and less discernible representations for nearby tokens. As a consequence, PI generally exhibits suboptimal results, especially as the extension ratio increases.

\textbf{NTK-based interpolation and extrapolation}. Prior studies such as~\cite{ntk,dynamicntk} analyze RoPE using an information encoding viewpoint and employ Neural Tangent Kernel (NTK) theory~\cite{jacot2018neural,tancik2020fourier} in their approach. To address the positional overlap problem inherent in PI, they propose redistributing interpolation effects among the various RoPE dimensions. Specifically, this involves applying less scaling to dimensions with higher frequency and more to those with lower frequency, thus facilitating both the interpolation and extrapolation of positions, where the scaling factor is $\lambda=s^{i}$. The dynamic NTK method~\cite{dynamicntk} further refines this by adapting the extension ratio for each position according to the current sequence length. In contrast to PI, which requires additional fine-tuning, NTK-based techniques are capable of lengthening the context window without the need for fine-tuning, though the maximum attainable extension is generally limited to 4$\times$.

\textbf{YaRN}~\cite{yarn} presents a notable advance in enhancing positional interpolation. The method categorizes RoPE dimensions into three distinct sets based on frequency characteristics, assigning a unique interpolation mechanism to each. Specifically, dimensions with high frequencies are subject to extrapolation ($\lambda$=1), whereas those with low frequencies utilize linear interpolation (PI). The intermediate frequency dimensions are handled using NTK. A central aspect of YaRN is its classification of RoPE dimensions, which is currently derived from manual, empirically-driven procedures. This reliance on human judgment may hinder optimal effectiveness when applied to different LLM architectures.

\subsection{Study on Non-uniform Positional Interpolation}

Drawing inspiration from NTK and YaRN, we observe that their improvements stem from leveraging non-linear approaches, particularly by treating various frequencies along RoPE dimensions differently to enhance both interpolation and extrapolation. Nonetheless, existing non-linear methods are predominantly built upon manually crafted heuristics. This observation leads us to consider two important questions: (1) Does the present method for positional interpolation represent the best possible approach? (2) Might there exist additional non-linear mechanisms that have yet to be discovered?

To address these questions, we employ evolution search (refer to Sec.~\ref{sec:method}) in order to identify improved non-uniform positional interpolations for LLaMA2-7B. Perplexity serves as the primary search criterion, assessed using five randomly selected samples from the PG19~\cite{pg19} validation dataset. Our empirical investigation yields the following principal insights.

\textbf{Finding 1}: \textit{RoPE dimensions exhibit substantial non-uniformities, which are not effectively handled by current positional interpolation methods.}

We optimize $\lambda$ individually for each RoPE dimension as specified in Eq.~\ref{eq:unified_rope}. Table~\ref{tbl:exp1} presents perplexity results for LLaMA2-7B across various approaches on the PG19 and Proof-pile~\cite{proof-pile} test sets, evaluated without fine-tuning. Our optimization yields substantial perplexity reductions, indicating that the commonly used linear (PI) and non-uniform (Dynamic-NTK and YaRN) interpolation methods do not achieve optimal performance. In particular, YaRN is found to perform worse than both PI and NTK on PG19, primarily because it fails to reach the intended context window when the LLM is not fine-tuned. For instance, with an 8k context length, perplexity with YaRN sharply increases after the 7k mark.

In our investigation, we observe that the rescaled factors $\lambda$ in Eq.~\ref{eq:unified_rope} are no longer uniform as in previous methods; they diverge from the constant scale $s$ utilized in PI, the NTK formula, and YaRN's group-based scaling. This departure from uniformity in the rescaling factors yields notable gains in language modeling performance (as measured by perplexity) for LLaMA2 across 8k and 16k context lengths, achieved without requiring fine-tuning. The performance improvement is attributable to the modified positional embedding, which more faithfully maintains the structure of the original RoPE—particularly across critical dimensions—thereby mitigating the model's challenge of differentiating among neighboring token positions.

\textbf{Finding 2}: \textit{RoPE for the initial tokens in the input sequence should be extrapolated with less interpolation.} 

We posit that for the first $\hat{n}$ tokens in the input sequence, reducing interpolation within their RoPE is desirable. The rationale lies in the disproportionately high attention scores these tokens acquire, making them integral to attention layers—a pattern previously noted in works such as Streaming LLM~\cite{streamingllm} and LM-Infinite~\cite{han2023lminfinite}. To test this assumption, we extended the context window to 8k and 16k tokens employing both PI and NTK approaches, but withheld interpolation for the initial $\hat n$ (0,2, ..., 256) tokens. Setting $\hat n=0$ returns us to the standard implementations of PI and NTK. As reported in Table~\ref{tbl:tokenposition}, two principal findings emerge: \textbf{(1)} omitting position interpolation for the earliest tokens yields measurable performance gains; \textbf{(2)} the choice of $\hat n$ that maximizes effectiveness is contingent on the length of the desired context extension.

\textbf{Finding 3}: \textit{Non-uniform positional interpolation effectively extends LLM context window in both fine-tuning and non-fine-tuning settings.} 

Our results indicate that although our discovered non-uniform position interpolation yields considerable enhancements in extension performance for 8k and 16k sequences without additional training, handling even longer contexts necessitates fine-tuning. Consequently, we apply fine-tuning to LLaMA2-7B using our optimized RoPE configuration for context windows up to 64k tokens (detailed settings are available in the Appendix). According to Table~\ref{tbl:finetune}, our approach consistently surpasses both PI and YaRN in performance, regardless of whether LLaMA2-7B has undergone fine-tuning. This superior performance can be attributed to our robust exploitation of non-uniform positional interpolation, which better preserves information and facilitates improved model initialization for subsequent fine-tuning.

\textbf{Summary}. This research highlights two distinct sources of non-uniformity: those found across RoPE dimensions and those among token positions. Leveraging these aspects within positional interpolation methods substantially advances large language model context extension capabilities.

\section{LongRoPE}

\label{sec:method}
Building upon these insights, we introduce LongRoPE. This approach begins by incorporating an effective search algorithm designed to leverage both types of non-uniformity, and subsequently applies this technique to expand the LLM context window past 2 million tokens.

\subsection{Problem Formulation}

These dual forms of non-uniformity result in an expansive solution landscape and add layers of difficulty to the optimization process. To manage this, we recast the challenge of optimizing multidimensional non-uniform position interpolation as a search-based problem.

Suppose we have a LLM configured for a context window of size $L'$ and receive long-form inputs $\mathbf{X}$, where every sequence $\mathbf{x}\in\mathbf{X}$ contains more tokens than $L'$. The canonical rotary angle associated with the $i^{th}$ RoPE embedding dimension at token index $n$ is given by $\frac{n}{\beta^i}$. The optimization task is thus expressed as:

\begin{equation}
\begin{gathered}
		\label{eq:problem}

\underset {\mathbf{x}\in\mathbf{X}; \, |\mathbf{x}|\ge L'}{ \operatorname {arg\,min} } \,  \mathcal{L}\, \left( \text{LLM} (\text{RoPE}, \mathbf{X})\right), \text{with}
\\
\scalebox{0.93}{
$\underset {\begin{subarray}{l} i=0,\cdot\cdot,\frac{d}{2}-1;\\ n\in[ 0,|\mathbf{x}|);\end{subarray}}{\text{RoPE}_(n)}= {\left[\cdot\cdot, \cos\left(\mathbb{I}(\hat{\lambda}_{i}, \hat{n})\cdot\frac{n}{\beta^i}\right), \sin\left(\mathbb{I}(\hat{\lambda}_{i}, \hat{n})\cdot\frac{n}{\beta^i}\right), \cdot\cdot\right] }$}\\
\text{here} \; \mathbb{I}(\hat{\lambda}_{i},\hat{n}) = \begin{cases}
1 &\mbox{ if $n < \hat{n}$}\\
{\frac{1}{\lambda}_{i}} &\mbox{ if $n \geq \hat{n}$}
\end{cases}
	\end{gathered}
\end{equation}

In this approach, we define a collection of rescale factors, $\mathbb{I}(\hat{\lambda}_{i},\hat{n})$, which serve to account for two different types of non-uniformities. Here, $\hat{\lambda_i}$ represents the degree of non-uniformity across RoPE dimensions, while $\hat{n}$ specifies the irregularity concerning token positions. In particular, the function $\mathbb{I}(\hat{\lambda}_{i},\hat{n})$ is applied to modify the rotation angle associated with the $i^{th}$ RoPE dimension. The parameter $\hat\lambda_{i}$ acts as the scaling factor, and $\hat{n}$ indicates a positional threshold for the tokens. For positions before $\hat{n}$ (i.e., the first $\hat{n}$-1 tokens), the rescale factor $\hat\lambda_{i}$ is not utilized, so the unaltered RoPE rotary angle $\frac{n}{\beta^i}$ remains. Once the token position reaches $n\ge\hat{n}$, the scaling factor begins to adjust the angle.

Assuming a target context window of length $L'$, the task is to determine the optimal set of rescale factors ($\mathbb{I}(\hat{\lambda}_{0},\hat{n})$, $\mathbb{I}(\hat{\lambda}_{1},\hat{n})$, ... , $\mathbb{I}(\hat{\lambda}_{i},\hat{n})$, ...) corresponding to each RoPE dimension from the first to the $d^{th}$. Employing these optimized rescale factors within the $\text{RoPE}$ mechanism enables the target $\text{LLM}$ to minimize the next token prediction loss, $\mathcal{L}$ (that is, the perplexity), when processing input sequences $\mathbf{X}$ whose token length equals $L'$.

\subsection{Searching the Non-uniform Position Interpolation}

To address the challenge outlined in Eq.~\ref{eq:problem}, we propose a straightforward yet robust approach that systematically searches for optimal RoPE rescale factors, aiming to leverage the inherent multidimensional non-uniformities present in position embedding.

\textbf{Search space}. We construct an extensive search space capable of accommodating both types of non-uniformity, as summarized in Table~\ref{tbl:searchspace}. In particular, our framework permits individual rescale factors for each dimension within the RoPE embedding. For streamlined search space formulation, we optimize over $\lambda_i$ and $\hat{n}$, rather than directly tuning $\mathbb{I}(\hat{\lambda}_{i},\hat{n})$, where $\hat{\lambda}_{i}=1/\lambda_i$. 
Table~\ref{tbl:searchspace} details this configuration: $\lambda_i$ is chosen from an interval starting at 1.0 (corresponding to direct extrapolation) up to $s\times1.25$ (which enables even broader interpolation than PI), with increments of 0.01, where $s$ denotes the factor by which the target context window is expanded.

The parameter $\hat{n}$ specifies how many of the earliest token positions maintain their original RoPE embedding, foregoing position interpolation. In practice, we select $\hat{n}$ from the set \{0, 1, 2, 4, 8, 12, 16, 20, 24, 28, 32, 64, 128, 256\}. If $\hat{n}=0$, then the searched rescale factors are applied across all token positions.

\textbf{Evolution-based search}. The search space detailed in Table~\ref{tbl:searchspace} encompasses a vast array of positional interpolation configurations, making comprehensive exploration computationally prohibitive. For instance, implementing an extension with $s=4\times$ results in $400^{128/2}\times14=4\times10^{167}$ possible selections. The dimensionality of the search space grows exponentially with increasing extension ratios. To efficiently navigate this immense landscape, we employ evolution search~\cite{guo2020single}, incorporating two optimization strategies that markedly enhance the effectiveness of the search. The overall process is presented in Algorithm~\ref{alg:search}.

\textit{Optimized initial population generation}. Rather than relying solely on random initialization for the population of $P$ rescale factors, we explicitly include the three RoPE rescale factors associated with PI, NTK, and YaRN as dedicated individuals in the starting population. For the other $P$-3 members, we generate them by applying random mutations to these three rescale factors, each with a mutation probability of $p$.

\textit{Monotonically non-decreasing constraint}. Once the initial population is established, we evaluate each individual by calculating its LLM perplexity. This involves implementing the designated RoPE rescale factors within the target LLM and assessing the perplexity on the input $\mathbf{X}$. The highest-performing $k$ individuals are then selected as parents for the evolutionary process. Nonetheless, because of the sheer breadth of the search space, simple mutation and crossover can often generate low-quality solutions, resulting in a large number of redundant perplexity evaluations. This inefficiency becomes more pronounced for large $L'$, as each perplexity assessment requires significant computational time.

To mitigate this issue, a monotonicity constraint is enforced on the sampled RoPE rescaling factors, ensuring that $\lambda_i\le \lambda_{i+1}$ for all $i$. Only those RoPE configurations that fulfill this non-decreasing requirement are utilized for evaluating perplexity in LLMs, which helps to greatly curtail the computational search space. Concretely, we stipulate that $\lambda_i$ must increase monotonically along the RoPE dimension indices ($i=0, \ldots, 63$). This approach is grounded in NTK theory~\cite{jacot2018neural,tancik2020fourier,ntk}, which indicates that higher-frequency, lower-index dimensions necessitate less interpolation (i.e., smaller $\lambda_i$ values), while lower-frequency, higher-index dimensions can tolerate more interpolation (i.e., larger $\lambda_i$ values).

\begin{algorithm}[t]
	\small
	\caption{The search algorithm}
	\textbf{Input:} target LLM, input samples $\mathbf{X}$,   population size $P$, mutation size $N_1$, crossover size $N_2$, max iterations $\mathcal{T}$, mutate probability $p$\\

	\begin{algorithmic}[1]
		\label{alg:search}
		\STATE Top-k $\leftarrow \phi$;	
		\STATE $\text{P}_0 \leftarrow \textit{Initialize\_population\_with\_optimization}(P, \mathbf{X}, p)$; 
		\FOR{$i$ from $1$ to $\mathcal{T}$}
		\STATE \textit{Compute\_perplexity}(\text{LLM}, $\text{P}_{i-1}$, $\mathbf{X}$);
		\STATE  Top-k $\leftarrow$ \textit{Update\_Topk}(Top-k, $\text{P}_{i-1}$);
		\STATE $\text{P}_{mutation}$ $\leftarrow$ \textit{Mutation\_with\_mono\_constraint}(Top-k, $N_1$, $p$); 
		\STATE $\text{P}_{crossover}$ $\leftarrow$ \textit{Crossover\_with\_mono\_constraint}(Top-k, $N_2$);
		\STATE $\text{P}_i = \text{P}_{mutation} \cup \text{P}_{crossover} \cup $ Top-k;
		\ENDFOR
		\STATE Return the member in Top-k exhibiting the minimal perplexity;
	\end{algorithmic}
\end{algorithm}

\textbf{8$\times$ extension without fine-tuning}. Utilizing evolutionary search, we successfully discover non-uniform RoPE rescale factors that prioritize critical dimensions and positional indices, thereby reducing information loss caused by interpolation. As illustrated in Fig.\ref{fig:non-ft}, our approach enables LLaMA2’s context window to be expanded from 4k up to 32k tokens in the absence of fine-tuning. By comparison, other methods like PI, as well as non-uniform NTK and YaRN, exhibit substantial increases in perplexity once the context is extended beyond 2$\times$.

\subsection{Extending LLM Context Window to 2048K}
\label{sec:extendto2048k}

\textbf{Progressive extension to 2048k}. Here, we propose a technique for expanding the context window of pre-trained LLMs beyond the standard 4k tokens, reaching up to 2048k. Our approach utilizing non-uniform positional interpolation enables up to an 8$\times$ extension without the need for fine-tuning. However, when targeting substantially larger expansions (such as 512$\times$), fine-tuning becomes essential. A possible strategy involves identifying optimal RoPE scaling factors suitable for the intended 2048k window and subsequently performing fine-tuning. Despite this, such a process is hindered by the significant computational resources required. Additionally, our empirical findings suggest that effective fine-tuning becomes increasingly difficult as the extension factor grows (refer to Appendix).

Fortunately, LongRoPE demonstrates efficacy with both the base and fine-tuned extended LLM variants. As such, we propose an efficient and incremental approach that reaches the intended 2048k length using only 1k fine-tuning iterations, while keeping the training context length capped at 256k.

$\Diamond$ \textit{LongRoPE search for scaling pre-trained LLMs up to 256k}. Using LLaMA2 as the reference model, we experiment with increasing the maximum context window to 128k and 256k tokens. In this process, the length extension factors applied are 32$\times$ and 64$\times$, respectively.

$\Diamond$ \textit{Fine-tuning to 256k}. Following pre-training, we adapt the LLM to handle a 256k context window through a two-stage fine-tuning process. Initially, LLaMA2 undergoes 400 fine-tuning steps with RoPE rescaling parameters set for 128k. After this, we modify the RoPE rescaling factors to correspond to 256k on the resulting checkpoint, and proceed with an extra 600 fine-tuning steps. This staged approach demonstrates greater efficiency compared to immediate fine-tuning at the 256k context length.

$\Diamond$ \textit{Expanding a fine-tuned, length-extended LLM to 2048k using LongRoPE search}. In the final step, we apply a subsequent search procedure to an LLM that has already been fine-tuned to support a context length of 256k. Through this method, we achieve a substantial enlargement of the context window, reaching 2048k tokens, without the need for additional fine-tuning. This corresponds to a context length increase by a factor of $512\times$.

\textbf{Recovery in shorter context windows}. When scaling up the context window to a massive 2048k, a degradation in performance is observed within the initial context window. This phenomenon is a recognized limitation of positional interpolation~\cite{pi}, where the embedding vectors for positions in the original context are compressed into a significantly smaller space in higher dimensions, leading to compromised language model performance. Specifically, using an extension ratio of 512$\times$, positional representations for the initial 4k window are heavily compressed, exacerbating the issue.

To address this issue, we conduct an additional evolution search specifically for the extended LLM, aiming to fine-tune the RoPE rescale parameters for relatively short context lengths such as 4k and 8k. Since shorter contexts necessitate minimal positional interpolation, we limit the upper bound for the searched $\lambda$. At inference time, the LLM adapts the relevant RoPE rescale factors on-the-fly.

\section{LongRoPE}

\label{sec:method}
Building on these insights, we propose LongRoPE, a method that leverages a novel, efficient search algorithm to take full advantage of the observed non-uniformities. Utilizing this approach, we succeed in extending the LLM context window past 2 million tokens.

\subsection{Problem Formulation}

These two forms of non-uniformity significantly expand the solution space and complicate the optimization process. To tackle these challenges, we recast the problem of multidimensional non-uniform position interpolation optimization as a search task.

Considering an LLM designed for a context window of size $L'$, and input documents $\mathbf{X}$, where each segment $\mathbf{x} \in \mathbf{X}$ has a token count exceeding $L'$, let us represent the initial rotary angle in the RoPE embedding for the $i^{th}$ dimension at token index $n$ by $\frac{n}{\beta^i}$. Accordingly, the associated optimization task can be formalized as:

\begin{equation}
\begin{gathered}
		\label{eq:problem}

\underset{\mathbf{x} \in \mathbf{X};\, |\mathbf{x}| \ge L'}{\operatorname{arg\,min}}\, \mathcal{L}\left(\text{LLM}\left(\text{RoPE}, \mathbf{X}\right)\right), \quad \text{where}
\\
\scalebox{0.93}{
$\underset{\begin{subarray}{l} i = 0, \ldots, \frac{d}{2}-1;\\ n \in [0, |\mathbf{x}|);\end{subarray}}{\text{RoPE}_\left(n\right)} = \left[\ldots, \cos\left(\mathbb{I}\left(\hat{\lambda}_{i}, \hat{n}\right)\frac{n}{\beta^i}\right),\; \sin\left(\mathbb{I}\left(\hat{\lambda}_{i}, \hat{n}\right)\frac{n}{\beta^i}\right), \ldots \right]$}
\\
\text{and} \quad \mathbb{I}(\hat{\lambda}_{i}, \hat{n}) = \begin{cases}
1 & \text{if } n < \hat{n} \\
\frac{1}{\lambda_i} & \text{if } n \geq \hat{n}
\end{cases}
\end{gathered}
\end{equation}

In this context, we define a collection of scaling factors, $\mathbb{I}(\hat{\lambda}_{i},\hat{n})$, aimed at addressing both types of non-uniformities present in the model. Here, $\hat{\lambda}_{i}$ and $\hat{n}$ refer to the non-uniform characteristics associated with the RoPE dimension indices and token position indices, respectively. In particular, $\mathbb{I}(\hat{\lambda}_{i},\hat{n})$ modifies the rotation angle corresponding to the $i^{th}$ RoPE dimension, employing $\hat\lambda_{i}$ as the dimension-specific scaling parameter and $\hat{n}$ as the critical token position threshold. For token positions less than $\hat{n}$, that is, for the first $\hat{n}-1$ tokens, the scaling factor $\hat\lambda_{i}$ does not apply, and the standard RoPE rotational angle $\frac{n}{\beta^i}$ is maintained. The adjustment via the rescale factor is introduced only for positions where $n\ge\hat{n}$.

Assuming a desired context window length of $L'$, the task is to identify the optimal rescale factors ($\mathbb{I}(\hat{\lambda}_{0},\hat{n})$, $\mathbb{I}(\hat{\lambda}_{1},\hat{n})$, ...$\mathbb{I}(\hat{\lambda}_{i},\hat{n})$...) across each RoPE dimension from the $1^{st}$ to $d^{th}$. Through this, the adapted $\text{LLM}$—incorporating the rescaled $\text{RoPE}$—should be capable of minimizing the next-token prediction loss $\mathcal{L}$ (i.e., perplexity) for given input sequences $\mathbf{X}$ of length $L'$.

\subsection{Searching the Non-uniform Position Interpolation}

To address the issue described in Eq.~\ref{eq:problem}, we present a straightforward yet robust technique that identifies optimal RoPE rescale factors, enabling effective utilization of multidimensional irregularities present in positional embeddings.

\textbf{Search space}. Our approach constructs an extensive search space that comprehensively incorporates both types of non-uniformity. The details of this search space are depicted in Table~\ref{tbl:searchspace}. More specifically, we permit the optimization of unique rescale factors for each dimension within the RoPE embedding. To streamline the formulation of the search space, the search is conducted over $\lambda_i$ and $\hat{n}$, rather than directly over $\mathbb{I}(\hat{\lambda}_{i},\hat{n})$, with the relationship $\hat{\lambda}_{i}=1/\lambda_i$. Table~\ref{tbl:searchspace} further details that $\lambda_i$ can be selected from a minimum of 1.0 (which equates to direct extrapolation) up to $s\times1.25$ (representing a wider range of interpolation than PI), using increments of 0.01, where $s$ corresponds to the ratio by which the context window is to be expanded.

The parameter $\hat{n}$ determines how many of the earliest token positions maintain the original RoPE embedding, foregoing position interpolation. In practice, $\hat{n}$ is tuned over the set \{0, 1, 2, 4, 8, 12, 16, 20, 24, 28, 32, 64, 128, 256\}. Setting $\hat{n}=0$ means that the searched rescale factors are applied across every token position.

\textbf{Evolution-based search}. The search space described in Table~\ref{tbl:searchspace} includes a vast array of possible solutions for positional interpolation, making effective exploration highly challenging. As an illustration, an extension factor of $s=4\times$ results in $400^{128/2}\times14=4\times10^{167}$ possible configurations. The search space size increases exponentially with greater extension ratios. To efficiently navigate this space, we employ evolution search~\cite{guo2020single} along with two additional optimization strategies that considerably enhance search effectiveness. The entire search process is outlined in Algorithm~\ref{alg:search}.

\textit{Optimized initial population generation}. Rather than assigning all $P$ rescale factors in the population at random, we explicitly incorporate the RoPE rescale factors for PI, NTK, and YaRN as three designated members of the initial population. The other $P$-3 individuals are then produced by applying random mutations to these three base rescale factors, each with a probability $p$.

\textit{Monotonically non-decreasing constraint}. Once the initial population is formed, we evaluate each individual's LLM perplexity. To do so, we apply the associated RoPE rescale factors within the target LLM and then determine the perplexity on the input $\mathbf{X}$. The individuals with the top-k perplexity scores are then selected to serve as parents for the evolutionary process. Nevertheless, given the enormous search space, unguided mutation and crossover operations may frequently yield suboptimal candidates, resulting in redundant perplexity evaluations. This inefficiency becomes more pronounced when $L'$ is large, due to the computational expense of each perplexity inference.

To tackle this, a monotonic non-decreasing constraint is enforced on the set of sampled RoPE rescaling factors: $\lambda_i\le \lambda_{i+1}$. Only those RoPE variants that adhere to this restriction are used in LLM perplexity assessments, which substantially curtails the computational overhead of the search. In detail, we enforce that $\lambda_i$ grows monotonically alongside the RoPE dimension ($i$=0,...,63). This requirement of monotonicity across dimensions is supported by NTK theory~\cite{jacot2018neural,tancik2020fourier,ntk}, which indicates that lower-indexed (higher-frequency) dimensions benefit from less interpolation (i.e., a smaller $\lambda_i$), whereas higher-indexed (lower-frequency) dimensions are able to accommodate more interpolation (i.e., a greater $\lambda_i$).

\begin{algorithm}[t]
	\small
	\caption{The search algorithm}
	\textbf{Input:} target LLM, input samples $\mathbf{X}$, population size $P$, mutation size $N_1$, crossover size $N_2$, maximum number of iterations $\mathcal{T}$, mutation probability $p$\\

	\begin{algorithmic}[1]
		\label{alg:search}
		\STATE Top-k $\leftarrow \phi$;	
		\STATE $\text{P}_0 \leftarrow$ \textit{Initialize\_population\_with\_optimization} ($P$, $\mathbf{X}$, $p$); 
		\FOR{$i=1$ to $\mathcal{T}$}
		\STATE \textit{Compute\_perplexity}(LLM, $\text{P}_{i-1}$, $\mathbf{X}$);
		\STATE Top-k $\leftarrow$ \textit{Update\_Topk}(Top-k, $\text{P}_{i-1}$);
		\STATE $\text{P}_{mutation} \leftarrow$ \textit{Mutation\_with\_mono\_constraint}(Top-k, $N_1$, $p$); 
		\STATE $\text{P}_{crossover} \leftarrow$ \textit{Crossover\_with\_mono\_constraint}(Top-k, $N_2$);
		\STATE $\text{P}_i \leftarrow \text{P}_{mutation} \cup \text{P}_{crossover} \cup \text{Top-k}$;
		\ENDFOR
		\STATE Output the individual in Top-k that achieves the minimal perplexity;
	\end{algorithmic}
\end{algorithm}

\textbf{8$\times$ extension without fine-tuning}. Through evolutionary search, our approach discovers optimal non-uniform RoPE rescale factors that maintain critical dimension and position information, thereby reducing information degradation from interpolation. As shown in Fig.\ref{fig:non-ft}, this enables our method to expand LLaMA2's context window from 4k up to 32k tokens, all without the need for fine-tuning. Conversely, prior approaches like PI, as well as non-uniform NTK and YaRN, exhibit a notable increase in perplexity when scaling beyond a 2$\times$ window extension.

\subsection{Extending LLM Context Window to 2048K}
\label{sec:extendto2048k}

\textbf{Progressive extension to 2048k}. In this section, we present our strategy for expanding the context window of pre-trained LLMs from the standard 4k up to over 2048k. Our approach leverages non-uniform positional interpolation, enabling up to 8$\times$ context extension without any fine-tuning. However, when seeking much larger expansions—such as extending by a factor of 512$\times$—fine-tuning becomes essential. One possible approach involves identifying appropriate RoPE rescaling factors for the desired 2048k window and then performing fine-tuning. Yet, this process is often constrained by the immense computational and resource requirements. Furthermore, our practical observations indicate that effectively fine-tuning LLMs for such substantial extension ratios poses significant challenges (see Appendix).

Fortunately, LongRoPE demonstrates efficacy with both the base model as well as its fine-tuned variants configured for extended sequence lengths. To this end, we propose a resource-efficient, incremental approach that reaches the desired 2048k context window, requiring only 1k fine-tuning steps and utilizing training sequences capped at 256k tokens.

$\Diamond$ \textit{Investigating 256k context extension for pre-trained LLMs using LongRoPE search}. Using LLaMA2 as a case study, we perform a search aimed at enlarging the context window to 128k and 256k tokens. At this point, the corresponding extension ratios are 32$\times$ and 64$\times$.

$\Diamond$ \textit{Fine-tuning to 256k}. Next, we adapt the pre-trained LLM to extend its context window to 256k tokens. In detail, we initiate fine-tuning of LLaMA2 for 400 steps employing the RoPE rescaled factors set for 128k. After this phase, we swap in the RoPE rescaled factors adjusted for 256k on the resulting checkpoint, followed by a further 600 fine-tuning steps. This staged approach demonstrates greater efficiency compared to fine-tuning to 256k in a single phase.

$\Diamond$ \textit{Scaling up a fine-tuned extended LLM to 2048k using LongRoPE search}. In the final stage, we carry out an additional search procedure on the LLM previously fine-tuned for sequences of length 256k. Through this approach, we achieve a substantial expansion of the context window, reaching 2048k tokens, all without any additional fine-tuning. This results in a total expansion factor of $512\times$.

\textbf{Restoration within a reduced context window}. Upon expanding the context window to an extensive 2048k length, we observe diminished performance inside the initial context window. This effect has previously been documented in relation to positional interpolation~\cite{pi}, since the position embeddings for tokens located in the original context window are now confined to a considerably smaller subspace in the embedding space, resulting in degraded language model capabilities. Specifically, applying a 512$\times$ context window expansion causes the embeddings associated with the original 4k context window to be especially densely packed.

To address this issue, we conduct an additional evolutionary search on the expanded LLM, specifically targeting the adjustment of RoPE rescale factors for shorter context sizes such as 4k and 8k. For these shorter lengths, the upper bound of the searched $\lambda$ is decreased, reflecting the reduced need for positional interpolation. At inference time, the LLM autonomously modifies the appropriate RoPE rescale factors as needed.

 \section{Experiments}

\textbf{Evaluation Tasks and Models}. Our experiments incorporate LongRoPE into LLaMA2-7B and Mistral-7B, assessing their effectiveness across three primary criteria: (1) measuring perplexity when these extended-context LLMs process lengthy textual inputs; (2) evaluating the passkey retrieval task, which tests the model’s capacity to accurately identify a specific passkey hidden within largely unrelated text; and (3) benchmarking performance using standard LLM evaluation suites, restricted to a 4096-token context window.

\textbf{Fine-tuning}. For LLaMA2, the model is fine-tuned with a learning rate of $2 \times 10^{-5}$ that decays linearly, utilizing a global batch size of 32. The fine-tuning process encompasses 400 steps on the Redpajama~\cite{redpajama} dataset, which is divided into 128k token segments marked with BOS and EOS tokens. Building upon the resulting checkpoint, an additional 600 steps are conducted to expand the context window to 256k tokens. The 128k context model training is performed with 8 A100 GPUs and a distributed training system~\cite{cube}, while the 256k model requires scaling up to 16 A100 GPUs. For Mistral, training is performed with a fixed learning rate of $1 \times 10^{-6}$ and a global batch size of 64. In alignment with the YaRN configuration~\cite{yarn}, both the 128k and 256k context models are trained for 400 steps on the Together Computer’s Long-Data Collections~\cite{mistraldata}, utilizing sequences of 16k tokens, with the training executed on 4 A100 GPUs.

\textbf{Search}. For a target window size up to 256k, the following configuration is applied: $P$=64, $N_1$=$N_2$=16, $p$=0.3, and $\mathcal{T}$=40, with the top-32 candidates selected for mutation and crossover in each cycle. Perplexity assessments are performed using five randomly chosen samples from the PG19 validation set, each meeting a minimum required context length. For target windows exceeding 512k, we reduce the population, as well as the mutation and crossover group sizes, by half. In these cases, perplexity is evaluated using three randomly selected examples from the Pile-Books3~\cite{pile} validation split.

\textbf{Baselines}. To achieve a 2048k context window, we conducted fine-tuning on models previously trained with 128k and 256k context lengths. The resulting models are named LongRoPE-2048k (ft=128k) for LLaMA2 and LongRoPE-2048k (ft=256k) for Mistral. Our evaluation involves a comparison among these four models and several leading methods for extending context windows, particularly focusing on open-source LLMs that have undergone fine-tuning post-positional interpolation via PI, NTK, and YaRN. Included in the baseline comparisons are Together-32k~\cite{together}, Code LLaMA~\cite{codellama}, LongLoRA-full-FT-100k~\cite{longlora}, as well as YaRN-LLaMA and YaRN-Mistral~\cite{yarn}.

\subsection{Main Results}

\label{sec:mainresult}
\textbf{Long sequence language modeling within 256k}. Our analysis first addresses how our approach compares to cutting-edge extended LLMs when evaluated with sequence lengths up to 256k. To showcase the general applicability of our method, we conduct experiments using two datasets: Proof-pile~\cite{pg19} and the test set from PG19~\cite{pile}. Perplexity scores are computed over a range of context sizes via a 256-token sliding window. For PG19, evaluation is performed on its entire test split, consisting of 100 documents. Following the setup of YaRN~\cite{yarn}, we sample 10 instances from Proof-pile, ensuring that each selected sample contains at least 128k tokens.

Table~\ref{tbl:proof-pile} and Table~\ref{tbl:pg19} present a comparison of the perplexity scores for LLaMA2 and Mistral when extended through various interpolation approaches on the Proof-pile and PG19 benchmarks, respectively. Two primary findings emerge: \textbf{(1)} The extended models demonstrate a consistent reduction in perplexity as the evaluation context increases from 4k up to 256k, indicating effective utilization of extended context windows. \textbf{(2)} Notably, even under the demanding scenario of a 16$\times$ expanded context window—which generally hinders models’ performance on shorter sequences—our LongRoPE-2048k configurations surpass state-of-the-art baselines within the 256k context range.

\textbf{Long sequence language modeling beyond 2000k}. To assess performance on exceptionally lengthy texts, we utilize the Books3~\cite{pile} dataset. For practical evaluation, we draw a random sample of 20 books, all with lengths over 2048k, and employ a sliding window approach with a size of 256k.

As detailed in Table~\ref{tbl:books3}, LongRoPE is able to expand the context window of both LLaMA2-7B and Mistral-7B to 2048k tokens, and achieves perplexity scores that are on par with, or outperform, baseline models when evaluated on shorter context windows ranging from 8k to 128k. A clear performance gap emerges between the 2048k variants of LLaMA2 and Mistral: while Mistral surpasses baseline results at shorter window sizes, its perplexity rises above 7 for context windows greater than 256k. LLaMA2, as anticipated, maintains lower perplexity as context increases, with slight upticks observed only at 1024k and 2048k. Furthermore, for LLaMA2, the LongRoPE-2048k variant exhibits superior performance when fine-tuned with a 256k window in comparison to 128k, attributed to the reduced secondary extension factor (8$\times$ versus 16$\times$). Conversely, Mistral attains better results with a 128k fine-tuning window. This discrepancy primarily arises because, in fine-tuning Mistral at 128k and 256k window sizes, the YaRN protocol is used, which limits the training length to 16k, thereby constraining Mistral’s capacity to further increase its effective context following fine-tuning.

\textbf{Passkey retrieval}. Next, we analyze the impact of context window length on generative tasks by employing a synthetic passkey retrieval benchmark as described in~\cite{passkey}. In this setup, a language model must identify a concealed passkey (specifically, a five-digit numeral) embedded within an extended input document. The format of the prompt can be found in the appendix. For evaluation, we conduct 10 rounds of the passkey retrieval task, randomly positioning the passkey at locations uniformly sampled from the available evaluation context span.

Fig.~\ref{fig:passkey} presents the comparison of retrieval accuracy against baseline models. The accuracy for existing methods falls sharply to zero when surpassing 128k tokens. In sharp contrast, LongRoPE-LLaMA2-2048k (ft=256k) demonstrates robust performance, sustaining a high retrieval accuracy ($\ge$90\%) across a wide range, from 4k up to 2048k, even though extracting a passkey from a corpus of millions of tokens is a difficult problem. For LongRoPE-Mistral-2048k (ft=128k), the model preserves perfect (100\%) accuracy up to 1800k tokens, then decreases to 60\% accuracy at 2048k. This outcome is consistent with the observations in Table~\ref{tbl:books3}, which reports a slight perplexity increase at the 2048k length.

\textbf{Standard benchmarks within the original context window}. To assess the performance of LongRoPE-2048k models over their native context window, we utilize the Hugging Face Open LLM Leaderboard~\cite{huggingfaceboard}, conducting both zero-shot and few-shot evaluations. Specifically, our benchmarks include 25-shot ARC-Challenge~\cite{allenai:arc}, 10-shot HellaSwag~\cite{zellers2019hellaswag}, 5-shot MMLU~\cite{mmlu}, and 0-shot TruthfulQA~\cite{lin2021truthfulqa}.

According to Table~\ref{tbl:commonsense}, our models produce results similar to those on the original benchmark, which was intended for use with a shorter context window. Notably, they surpass the original Mistral on TruthfulQA with an improvement of +0.5\%. Although LongRoPE-LLaMA2-2048k, which underwent fine-tuning at 256k, demonstrates a marginally greater drop in performance, its effectiveness largely stays within acceptable limits across the majority of tasks.

\subsection{Ablation Results}

\textbf{Effectiveness of the second positional interpolation}. Within the framework of our progressive extension methodology, we implement a second stage of non-uniform positional interpolation through our search algorithm applied to fine-tuned extended LLMs. To assess its performance, we experiment with the LLaMA2-256k model that has been fine-tuned. The model is subsequently extended to 512k, 1024k, and 2048k contexts using both PI and YaRN methods. As indicated in Table~\ref{tbl:secondsearch}, our approach involving non-uniform positional interpolation maintains a stable perplexity across these larger context lengths. On the other hand, the perplexity of models extended with standard PI or YaRN rises significantly as the extension ratio increases.

\textbf{Recovery performance at reduced context lengths.} In order to address the decline in performance seen at lower context lengths, we fine-tune the RoPE factors for LongRoPE-2048k using our optimization algorithm. Concretely, the upper bound for scale factors during the search is reduced, aiming to limit interpolation for shorter context windows of 4k and 8k. Table~\ref{tbl:shortperformance} presents a comparison of perplexity for LongRoPE-LLaMA2-2048k on Proof-pile at both 4k and 8k tokens, alongside mean LLM benchmark accuracy. These outcomes reveal that considerable gains in performance are achieved when operating with shorter context spans.

\textbf{Analysis on the two forms of non-uniformities}. To assess the impact of each type of non-uniformity on model performance, we conduct an ablation study. Two experimental configurations are considered: (i) scaling LLaMA2-7B to sequence lengths of 16k and 32k with various strategies—PI, optimizing only the RoPE dimension, and jointly searching for both non-uniformities; (ii) adapting our 256k-length fine-tuned LLaMA2 model to a length of 2048k using analogous techniques. In all cases, perplexity is measured without additional fine-tuning. Results in Table~\ref{tbl:searchdimension} indicate that accounting for non-uniformity in the RoPE dimension notably decreases perplexity, outperforming PI’s linear interpolation. Adjusting for token position non-uniformity yields clear performance improvements at 16k and 32k lengths, but this effect diminishes at the extreme 2048k length, likely due to the challenging scale. Simply retaining initial tokens without interpolation is insufficient in these cases; investigating solutions for this regime remains an open area for future research.

\section{Related Works}

In addition to methods based on position interpolation, this section discusses related works of other approaches. 

\textbf{Retrieval-based} methods rely on an external memory component to store extended historical context and utilize retrieval mechanisms to fetch pertinent documents during inference~\cite{longllama,wang2023augmenting,borgeaud2022improving}. Implementing these strategies often requires substantial, explicit adjustments to the underlying LLM architectures. In comparison, our approach introduces only minimal changes to positional embeddings, offering a more streamlined solution. Additionally, we extend support to a wider range of long-context applications, including long document summarization and few-shot learning, going beyond traditional retrieval-centric tasks.

\textbf{Attention-based context window extensions}. In addition to positional embedding interpolation, several studies extend the LLM's context window by altering attention mechanisms, while retaining the model's original context window size~\cite{han2023lminfinite, streamingllm,parallelwindow}. The central principle involves the use of innovative attention masks designed to address the problem of increased attention complexity associated with additional positions. These approaches are complementary to positional interpolation strategies.

\textbf{Fine-tuning based approaches} concentrate on optimizing the adaptation of pre-trained LLMs to longer contexts by modifying position embeddings. Approaches such as Code LLaMA~\cite{codellama}, LLaMA2 Long~\cite{xiong2023effective}, and ScaledRoPE~\cite{liu2023scaling} implement this by selecting a substantially increased RoPE base value and then fine-tuning models to the intended sequence length. In contrast, our approach is designed to accommodate a broad spectrum of target lengths and is capable of extending sequence length to beyond 2M tokens. In response to the intensive GPU requirements associated with fine-tuning at large context sizes (e.g., above 128k tokens), methods like LongLoRA~\cite{longlora} and PoSE~\cite{zhu2023pose} have been introduced to reduce computational burden. Our method is complementary to these resource-efficient fine-tuning strategies.

\section{Conclusion}

In this paper, we introduce LongRoPE, a technique capable of expanding the context window of large language models (LLMs) to an unprecedented 2048k tokens, all while preserving their proficiency on tasks involving the original, shorter context lengths. Our method utilizes an efficient evolutionary search to exploit two types of non-uniformity within RoPE positional embeddings. This approach yields dual advantages: it delivers favorable initialization for fine-tuning and achieves up to an 8$\times$ increase in context length without necessitating fine-tuning. Building further, we employ a staged extension mechanism by fine-tuning LLMs at 256k context lengths and subsequently attaining a 2048k window without additional fine-tuning. Comprehensive experimental results confirm the efficacy of LongRoPE. We anticipate that LongRoPE-2048k models will unlock a wide range of new long-context use cases and foster subsequent advancements in the area.

\section*{Broader Impacts} The research described in this paper seeks to further the development of Machine Learning as a discipline. While our findings may have diverse implications for society, we have not identified any particular outcomes that warrant explicit mention at this time.

	\balance

\end{document}